\documentclass{article}
\usepackage[utf8]{inputenc}
\usepackage{cite}

\title{Data recovery in 5G wireless channels with doppler and delay effect using mixture of experts for enhanced Industry 4.0 mobility.}
%Recuperación de datos en canales inalámbricos 5G con efecto Doppler y retardo utilizando mezcla de expertos para mejorar la movilidad en la Industria 4.0
\author{Luis Emilio Tonix Gleason}
\date{\today}

\begin{document}
La Industria 4.0 representa un gran potencial para el futuro, como lo demuestra el crecimiento significativo de las inversiones en tecnologías relacionadas. Tan solo en 2019, las inversiones en IoT en India alcanzaron aproximadamente 5,000 millones de dólares estadounidenses \cite{iosr2021iot}. Estas tecnologías permiten recopilar, procesar y transferir datos en tiempo real, aplicándose a sectores críticos como la gestión energética e hidráulica en países desarrollados \cite{Water40}.

El crecimiento del Internet de las Cosas (IoT) está impulsado por la creciente expansión de dispositivos conectados y la adopción de tecnologías avanzadas como 4G, 5G y Wi-Fi, que dependen de la modulación OFDM y la codificación PSK para garantizar la eficiencia en las comunicaciones \cite{10729310}. Estas tecnologías son utilizadas diariamente en aplicaciones móviles y de oficina, pero también presentan desafíos significativos debido a condiciones ambientales adversas como altas temperaturas, humedad, acumulación de polvo, y la limitada capacidad de mantenimiento y personal capacitado \cite{8519869}.

Un aspecto menos explorado del IoT es su implementación en dispositivos de alta movilidad, donde los equipos operan a velocidades variables \cite{10152108}. Esta variabilidad genera distorsiones en la señal portadora, con efectos como el Doppler ambiguo, que comienza a ser significativo desde velocidades de 25 m/s y distancias de 500 m, alcanzando un mínimo de 5 kHz \cite{8022917}. Estas distorsiones afectan la precisión de los datos codificados y transmitidos, lo que plantea retos para aplicaciones que requieren una alta fiabilidad.

\bibliographystyle{plain}
\bibliography{references}


\end{document}